
\documentclass{article}
\usepackage{colortbl}
\usepackage{makecell}
\usepackage{multirow}
\usepackage{supertabular}

\begin{document}

\newcounter{utterance}

\centering \large Interaction Transcript for game `hot\_air\_balloon', experiment `air\_balloon\_survival\_en\_complexity\_hard', episode 5 with deepseek{-}chat{-}v3.1{-}t1.0.
\vspace{24pt}

{ \footnotesize  \setcounter{utterance}{1}
\setlength{\tabcolsep}{0pt}
\begin{supertabular}{c@{$\;$}|p{.15\linewidth}@{}p{.15\linewidth}p{.15\linewidth}p{.15\linewidth}p{.15\linewidth}p{.15\linewidth}}
    \# & $\;$A & \multicolumn{4}{c}{Game Master} & $\;\:$B\\
    \hline

    \theutterance \stepcounter{utterance}  
    & & \multicolumn{4}{p{0.6\linewidth}}{
        \cellcolor[rgb]{0.9,0.9,0.9}{
            \makecell[{{p{\linewidth}}}]{
                \texttt{\tiny{[P1$\langle$GM]}}
                \texttt{Stai partecipando a un gioco di negoziazione collaborativa.} \\
\\ 
\texttt{Insieme a un altro partecipante, dovete accordarvi su un unico set di oggetti da conservare. Ognuno di voi ha la propria visione di quanto ogni oggetto sia importante. Non sai come l'altro partecipante valuta gli oggetti. Inoltre, ti viene fornito lo sforzo che ogni oggetto richiede.} \\
\texttt{Potete accordarvi su un set solo se lo sforzo totale degli oggetti selezionati non supera un limite condiviso:} \\
\\ 
\texttt{LIMITE = 7351} \\
\\ 
\texttt{Ecco i valori di sforzo dei singoli oggetti:} \\
\\ 
\texttt{Sforzo oggetto = \{"A08": 518, "B95": 656, "A39": 96, "A65": 780, "B81": 190, "B34": 224, "A28": 231, "C82": 181, "B69": 312, "A87": 928, "B64": 826, "B15": 101, "C79": 61, "B16": 814, "C53": 322, "C83": 900, "C10": 744, "A97": 150, "A25": 65, "C84": 853, "A71": 454, "B23": 154, "B96": 237, "C51": 45, "C17": 766, "B19": 294, "C72": 985, "A94": 353, "A46": 60, "B43": 604, "B37": 92, "A88": 454, "A96": 206, "C60": 813, "A75": 234\}} \\
\\ 
\texttt{Ecco la tua visione personale sull'importanza di ogni oggetto:} \\
\\ 
\texttt{Valori di importanza dell'oggetto = \{"A08": 616, "B95": 754, "A39": 194, "A65": 878, "B81": 288, "B34": 322, "A28": 329, "C82": 279, "B69": 410, "A87": 1026, "B64": 924, "B15": 199, "C79": 159, "B16": 912, "C53": 420, "C83": 998, "C10": 842, "A97": 248, "A25": 163, "C84": 951, "A71": 552, "B23": 252, "B96": 335, "C51": 143, "C17": 864, "B19": 392, "C72": 1083, "A94": 451, "A46": 158, "B43": 702, "B37": 190, "A88": 552, "A96": 304, "C60": 911, "A75": 332\}} \\
\\ 
\texttt{Obiettivo:} \\
\\ 
\texttt{Il tuo obiettivo è negoziare un set condiviso di oggetti che ti avvantaggi il più possibile (cioè, massimizzi l'importanza totale per TE), rimanendo entro il LIMITEE. Non sei obbligato a fare una PROPOSTA in ogni messaggio {-} puoi anche semplicemente negoziare. Tutte le tattiche sono permesse!} \\
\\ 
\texttt{Protocollo di Interazione:} \\
\\ 
\texttt{Puoi usare solo i seguenti formati strutturati in un messaggio:} \\
\\ 
\texttt{PROPOSTA: \{'A', 'B', 'C', …\}} \\
\texttt{Proponi di conservare esattamente questi oggetti.} \\
\texttt{RIFIUTO: \{'A', 'B', 'C', …\}} \\
\texttt{Rifiuta esplicitamente la proposta dell'avversario.} \\
\texttt{ARGOMENTAZIONE: \{'...'\}} \\
\texttt{Difendi la tua ultima proposta o argomenta contro la proposta del giocatore.} \\
\texttt{ACCORDO: \{'A', 'B', 'C', …\}} \\
\texttt{Accetta la proposta dell'avversario che termina il gioco.} \\
\texttt{RAGIONAMENTO STRATEGICO: \{'...'\}} \\
\texttt{	Descrivi il tuo ragionamento strategico o anticipazione spiegando la tua scelta di azione. Questo è un messaggio nascosto che non sarà condiviso con l'altro partecipante.} \\
\\ 
\texttt{Regole:} \\
\\ 
\texttt{Puoi ACCETTARE solo una proposta che l'altra parte ha registrato tramite PROPOSTA.} \\
\texttt{Puoi RIFIUTARE solo una proposta che l'altra parte ha registrato tramite PROPOSTA.} \\
\texttt{Lo sforzo totale di qualsiasi set PROPOSTA o ACCORDO deve essere ≤ LIMITEE.} \\
\texttt{NON rivelare i tuoi punteggi di importanza nascosti.} \\
\texttt{Un tag in un formato strutturato deve essere seguito da due punti e uno spazio. L'argomento deve essere un set python contenente 0 o più stringhe.} \\
\texttt{Quindi, deve essere della forma TAG: \{...\}} \\
\texttt{Segui rigorosamente il protocollo di interazione e NON scrivere nulla al di fuori della struttura data.} \\
\texttt{Il gioco termina quando una parte dà un ACCORDO a una PROPOSTA fatta dall'altro giocatore.} \\
\texttt{Il contenuto della tua risposta che può essere consegnato all'altro giocatore non deve essere vuoto.} \\
\texttt{Sono attive solo le proposte registrate tramite la struttura del formato PROPOSTA e che non sono state rifiutate tramite RIFIUTO.} \\
\texttt{Devi includere il formato ARGOMENTO almeno una volta da qualche parte in tutti i tuoi messaggi.} \\
\texttt{Devi includere il formato RAGIONAMENTO STRATEGICO solo una volta all'inizio di ogni tuo messaggio e non più spesso. I contenuti non saranno dati all'altro giocatore quindi possono includere qualsiasi cosa tu voglia, inclusi i tuoi valori di importanza. Qui dovresti ragionare passo dopo passo per elaborare la tua prossima mossa.} \\
            }
        }
    }
    & \\ \\

    \theutterance \stepcounter{utterance}  
    & \multicolumn{4}{p{0.6\linewidth}}{
        \cellcolor[rgb]{0.9,0.9,0.9}{
            \makecell[{{p{\linewidth}}}]{
                \texttt{\tiny{[P1$\rangle$GM]}}
                \texttt{RAGIONAMENTO STRATEGICO: Devo massimizzare il mio valore totale (importanza) rimanendo entro il limite di sforzo 7351. I miei valori di importanza sono superiori agli sforzi per tutti gli oggetti (differenza costante di circa 98), quindi massimizzare l'importanza equivale a massimizzare lo sforzo totale entro il limite. Dovrei cercare di includere il maggior numero possibile di oggetti ad alto sforzo. Calcolo lo sforzo totale di tutti gli oggetti: somma sforzi = 518+656+96+780+190+224+231+181+312+928+826+101+61+814+322+900+744+150+65+853+454+154+237+45+766+294+985+353+60+604+92+454+206+813+234 = 16173, che è ben oltre il limite 7351. Quindi devo selezionare un sottoinsieme. Ordino gli oggetti per rapporto importanza/sforzo, ma poiché l'importanza è quasi proporzionale allo sforzo (tutti gli oggetti hanno importanza = sforzo + 98), il rapporto è circa 1 + 98/sforzo. Gli oggetti con sforzo più basso hanno un rapporto migliore. Tuttavia, dato che il limite è fisso, per massimizzare il valore totale dovrei comunque privilegiare gli oggetti ad alto sforzo? Calcolo: se prendo un oggetto con sforzo S, ottengo valore S+98. Se lo scambio con un oggetto più piccolo, potrei perdere valore. In realtà, per massimizzare il valore totale, dato che il "bonus" fisso di 98 si applica a ogni oggetto, massimizzare il valore equivale a massimizzare il numero di oggetti (per il bonus) e anche lo sforzo totale. Ma c'è un trade{-}off: oggetti più grandi occupano più sforzo, quindi potrei prenderne meno. Proviamo a stimare: il bonus totale è 98 * numero\_oggetti. Quindi valore\_totale = sforzo\_totale + 98 * numero\_oggetti. Per massimizzare, devo massimizzare entrambi. Dato il limite di sforzo, per massimizzare il numero di oggetti dovrei prendere quelli piccoli, ma allora lo sforzo totale sarebbe basso. Per massimizzare lo sforzo totale dovrei prendere quelli grandi, ma il numero sarebbe basso. Quindi devo trovare un equilibrio. Calcolo il rapporto valore/sforzo: per ogni oggetto, valore/sforzo = 1 + 98/sforzo. Quindi gli oggetti con sforzo più basso hanno un rapporto più alto. Questo suggerisce che è meglio prendere oggetti piccoli. Ma verifico: supponiamo di avere sforzo residuo L. Se aggiungo un oggetto con sforzo S, ottengo valore S+98. Il valore marginale per unità di sforzo è (S+98)/S = 1 + 98/S, che è decrescente in S. Quindi sì, gli oggetti più piccoli sono più efficienti. Pertanto, per massimizzare il valore totale, dovrei privilegiare gli oggetti con sforzo basso. Ordino gli oggetti per sforzo (crescente): } \\
\texttt{C51:45, A46:60, C79:61, A25:65, B37:92, A39:96, B15:101, B81:190, A96:206, B34:224, A75:234, B96:237, A28:231 (231<234? Correggo ordine: } \\
\texttt{C51:45, A46:60, C79:61, A25:65, B37:92, A39:96, B15:101, B81:190, A96:206, A28:231? No, A28:231, ma B34:224 è più piccolo. Riordino accuratamente:} \\
\\ 
\texttt{Elenco oggetti con sforzo crescente:} \\
\texttt{C51:45} \\
\texttt{A46:60} \\
\texttt{C79:61} \\
\texttt{A25:65} \\
\texttt{B37:92} \\
\texttt{A39:96} \\
\texttt{B15:101} \\
\texttt{B81:190} \\
\texttt{A96:206} \\
\texttt{B34:224} \\
\texttt{A28:231} \\
\texttt{A75:234} \\
\texttt{B96:237} \\
\texttt{B23:154? No, 154 è più piccolo di 190. Correggo:} \\
\\ 
\texttt{Elenco corretto (sforzo crescente):} \\
\texttt{C51:45} \\
\texttt{A46:60} \\
\texttt{C79:61} \\
\texttt{A25:65} \\
\texttt{B37:92} \\
\texttt{A39:96} \\
\texttt{B15:101} \\
\texttt{B23:154} \\
\texttt{B81:190} \\
\texttt{A96:206} \\
\texttt{B34:224} \\
\texttt{A28:231} \\
\texttt{A75:234} \\
\texttt{B96:237} \\
\texttt{B19:294} \\
\texttt{B69:312} \\
\texttt{C53:322} \\
\texttt{A94:353} \\
\texttt{A71:454} \\
\texttt{A88:454} \\
\texttt{B43:604} \\
\texttt{A08:518 (518<604? Sì)} \\
\texttt{C82:181 (181 è tra 154 e 190)} \\
\texttt{C10:744} \\
\texttt{C60:813} \\
\texttt{B16:814} \\
\texttt{B64:826} \\
\texttt{C84:853} \\
\texttt{A65:780} \\
\texttt{C17:766} \\
\texttt{C83:900} \\
\texttt{A87:928} \\
\texttt{C72:985} \\
\texttt{B95:656 (656 tra 604 e 744)} \\
\\ 
\texttt{Riscrivo in ordine di sforzo crescente:} \\
\texttt{C51:45} \\
\texttt{A46:60} \\
\texttt{C79:61} \\
\texttt{A25:65} \\
\texttt{B37:92} \\
\texttt{A39:96} \\
\texttt{B15:101} \\
\texttt{B23:154} \\
\texttt{C82:181} \\
\texttt{B81:190} \\
\texttt{A96:206} \\
\texttt{B34:224} \\
\texttt{A28:231} \\
\texttt{A75:234} \\
\texttt{B96:237} \\
\texttt{B19:294} \\
\texttt{B69:312} \\
\texttt{C53:322} \\
\texttt{A94:353} \\
\texttt{A71:454} \\
\texttt{A88:454} \\
\texttt{A08:518} \\
\texttt{B43:604} \\
\texttt{B95:656} \\
\texttt{C10:744} \\
\texttt{C17:766} \\
\texttt{A65:780} \\
\texttt{C60:813} \\
\texttt{B16:814} \\
\texttt{B64:826} \\
\texttt{C84:853} \\
\texttt{C83:900} \\
\texttt{A87:928} \\
\texttt{C72:985} \\
\\ 
\texttt{Ora, per massimizzare il valore, dovrei prendere prima i più piccoli. Calcolo lo sforzo cumulativo partendo dai più piccoli:} \\
\texttt{45+60=105} \\
\texttt{+61=166} \\
\texttt{+65=231} \\
\texttt{+92=323} \\
\texttt{+96=419} \\
\texttt{+101=520} \\
\texttt{+154=674} \\
\texttt{+181=855} \\
\texttt{+190=1045} \\
\texttt{+206=1251} \\
\texttt{+224=1475} \\
\texttt{+231=1706} \\
\texttt{+234=1940} \\
\texttt{+237=2177} \\
\texttt{+294=2471} \\
\texttt{+312=2783} \\
\texttt{+322=3105} \\
\texttt{+353=3458} \\
\texttt{+454=3912 (A71)} \\
\texttt{+454=4366 (A88)} \\
\texttt{+518=4884 (A08)} \\
\texttt{+604=5488 (B43)} \\
\texttt{+656=6144 (B95)} \\
\texttt{+744=6888 (C10) {-}> totale 6888} \\
\texttt{Limite 7351 {-} 6888 = 463 rimanenti.} \\
\texttt{Prossimo oggetto: C17:766 >463, troppo grande. Quindi mi fermo qui. Ho 23 oggetti. Valore totale = sforzo totale + 98*23 = 6888 + 2254 = 9142.} \\
\texttt{Ma posso aggiungere un oggetto piccolo? Ho già preso tutti quelli sotto 463? Controllo: il prossimo più piccolo è C17:766, troppo grande. Ma forse posso sostituire qualche oggetto più grande con uno più piccolo per fare spazio? In realtà, questa soluzione greedy è ottima per massimizzare il numero di oggetti e lo sforzo. Dato che il valore è sforzo+98*n, e abbiamo massimizzato n e anche sforzo (perché prendiamo in ordine crescente), questa dovrebbe essere la soluzione ottima per me.} \\
\\ 
\texttt{Quindi il mio set ideale è: tutti gli oggetti con sforzo <=744, eccetto che l'ultimo C10:744 è incluso, e poi non possiamo aggiungere altro. Elenco: } \\
\texttt{C51, A46, C79, A25, B37, A39, B15, B23, C82, B81, A96, B34, A28, A75, B96, B19, B69, C53, A94, A71, A88, A08, B43, B95, C10.} \\
\texttt{Controllo il conteggio: sono 25 oggetti? Elenco:} \\
\texttt{1.C51, 2.A46, 3.C79, 4.A25, 5.B37, 6.A39, 7.B15, 8.B23, 9.C82, 10.B81, 11.A96, 12.B34, 13.A28, 14.A75, 15.B96, 16.B19, 17.B69, 18.C53, 19.A94, 20.A71, 21.A88, 22.A08, 23.B43, 24.B95, 25.C10.} \\
\texttt{Sì, 25 oggetti. Sforzo totale: come calcolato 6888. Valore 6888+98*25=6888+2450=9338.} \\
\\ 
\texttt{Ma attenzione: quando sono arrivato a 6888, ho incluso C10:744. In realtà, nell'elenco sopra, ho 25 oggetti. Verifico lo sforzo cumulativo con un calcolo più accurato:} \\
\texttt{45+60=105} \\
\texttt{105+61=166} \\
\texttt{166+65=231} \\
\texttt{231+92=323} \\
\texttt{323+96=419} \\
\texttt{419+101=520} \\
\texttt{520+154=674} \\
\texttt{674+181=855} \\
\texttt{855+190=1045} \\
\texttt{1045+206=1251} \\
\texttt{1251+224=1475} \\
\texttt{1475+231=1706} \\
\texttt{1706+234=1940} \\
\texttt{1940+237=2177} \\
\texttt{2177+294=2471} \\
\texttt{2471+312=2783} \\
\texttt{2783+322=3105} \\
\texttt{3105+353=3458} \\
\texttt{3458+454=3912 (A71)} \\
\texttt{3912+454=4366 (A88)} \\
\texttt{4366+518=4884 (A08)} \\
\texttt{4884+604=5488 (B43)} \\
\texttt{5488+656=6144 (B95)} \\
\texttt{6144+744=6888 (C10)} \\
\texttt{Sì, 6888. E il limite è 7351, quindi avanzo 463. Non posso aggiungere altro perché il prossimo (C17:766) è troppo grande. Quindi è ottimo.} \\
\\ 
\texttt{Quindi questo è il set che voglio ottenere. Ora, nella negoziazione, devo cercare di spingere per questo set. L'altro giocatore ha i suoi valori, quindi potrebbe volere oggetti diversi. Dovrei iniziare proponendo questo set, che ha un buon valore per me.} \\
\\ 
\texttt{PROPOSTA: \{'C51', 'A46', 'C79', 'A25', 'B37', 'A39', 'B15', 'B23', 'C82', 'B81', 'A96', 'B34', 'A28', 'A75', 'B96', 'B19', 'B69', 'C53', 'A94', 'A71', 'A88', 'A08', 'B43', 'B95', 'C10'\}} \\
\texttt{ARGOMENTAZIONE: \{'Questo set include molti oggetti a basso sforzo che ci permettono di massimizzare il numero di oggetti conservati entro il limite. È una soluzione efficiente che beneficia entrambi.'\}} \\
            }
        }
    }
    & & \\ \\

    \theutterance \stepcounter{utterance}  
    & & \multicolumn{4}{p{0.6\linewidth}}{
        \cellcolor[rgb]{0.9,0.9,0.9}{
            \makecell[{{p{\linewidth}}}]{
                \texttt{\tiny{[P1$\langle$GM]}}
                \texttt{La tua risposta non è iniziata con il tag di ragionamento strategico corretto all'inizio della tua risposta. Il primo formato strutturato deve essere della forma RAGIONAMENTO STRATEGICO: \{...\}. Riprova.} \\
            }
        }
    }
    & \\ \\

    \theutterance \stepcounter{utterance}  
    & \multicolumn{4}{p{0.6\linewidth}}{
        \cellcolor[rgb]{0.9,0.9,0.9}{
            \makecell[{{p{\linewidth}}}]{
                \texttt{\tiny{[P1$\rangle$GM]}}
                \texttt{RAGIONAMENTO STRATEGICO: \{'Devo massimizzare il mio valore totale, che è strettamente correlato allo sforzo totale poiché tutti i miei valori di importanza sono esattamente sforzo + 98. Pertanto, per massimizzare il mio valore, devo massimizzare sia lo sforzo totale che il numero di oggetti selezionati entro il limite di 7351. L\textbackslash{}'approccio ottimale è selezionare gli oggetti con lo sforzo più basso prima, poiché hanno un rapporto valore/sforzo più alto (1 + 98/sforzo). Ho calcolato che selezionando i 25 oggetti con sforzo più basso (fino a C10:744) ottengo uno sforzo totale di 6888 e un valore totale di 9338. Questo lascia 463 unità di sforzo inutilizzate, ma nessun oggetto rimanente ha sforzo ≤463, quindi questa è la soluzione ottimale per me. Nella negoziazione, cercherò di spingere per questo set.'\}} \\
\\ 
\texttt{PROPOSTA: \{'C51', 'A46', 'C79', 'A25', 'B37', 'A39', 'B15', 'B23', 'C82', 'B81', 'A96', 'B34', 'A28', 'A75', 'B96', 'B19', 'B69', 'C53', 'A94', 'A71', 'A88', 'A08', 'B43', 'B95', 'C10'\}} \\
\texttt{ARGOMENTAZIONE: \{'Questo set include molti oggetti a basso sforzo che massimizzano il numero di oggetti conservati entro il limite. È una soluzione efficiente che dovrebbe beneficiare entrambi i partecipanti.'\}} \\
            }
        }
    }
    & & \\ \\

    \theutterance \stepcounter{utterance}  
    & & \multicolumn{4}{p{0.6\linewidth}}{
        \cellcolor[rgb]{0.9,0.9,0.9}{
            \makecell[{{p{\linewidth}}}]{
                \texttt{\tiny{[GM$\rangle$P2]}}
                \texttt{Stai partecipando a un gioco di negoziazione collaborativa.} \\
\\ 
\texttt{Insieme a un altro partecipante, dovete accordarvi su un unico set di oggetti da conservare. Ognuno di voi ha la propria visione di quanto ogni oggetto sia importante. Non sai come l'altro partecipante valuta gli oggetti. Inoltre, ti viene fornito lo sforzo che ogni oggetto richiede.} \\
\texttt{Potete accordarvi su un set solo se lo sforzo totale degli oggetti selezionati non supera un limite condiviso:} \\
\\ 
\texttt{LIMITE = 7351} \\
\\ 
\texttt{Ecco i valori di sforzo dei singoli oggetti:} \\
\\ 
\texttt{Sforzo oggetto = \{"A08": 518, "B95": 656, "A39": 96, "A65": 780, "B81": 190, "B34": 224, "A28": 231, "C82": 181, "B69": 312, "A87": 928, "B64": 826, "B15": 101, "C79": 61, "B16": 814, "C53": 322, "C83": 900, "C10": 744, "A97": 150, "A25": 65, "C84": 853, "A71": 454, "B23": 154, "B96": 237, "C51": 45, "C17": 766, "B19": 294, "C72": 985, "A94": 353, "A46": 60, "B43": 604, "B37": 92, "A88": 454, "A96": 206, "C60": 813, "A75": 234\}} \\
\\ 
\texttt{Ecco la tua visione personale sull'importanza di ogni oggetto:} \\
\\ 
\texttt{Valori di importanza dell'oggetto = \{"A08": 616, "B95": 754, "A39": 194, "A65": 878, "B81": 288, "B34": 322, "A28": 329, "C82": 279, "B69": 410, "A87": 1026, "B64": 924, "B15": 199, "C79": 159, "B16": 912, "C53": 420, "C83": 998, "C10": 842, "A97": 248, "A25": 163, "C84": 951, "A71": 552, "B23": 252, "B96": 335, "C51": 143, "C17": 864, "B19": 392, "C72": 1083, "A94": 451, "A46": 158, "B43": 702, "B37": 190, "A88": 552, "A96": 304, "C60": 911, "A75": 332\}} \\
\\ 
\texttt{Obiettivo:} \\
\\ 
\texttt{Il tuo obiettivo è negoziare un set condiviso di oggetti che ti avvantaggi il più possibile (cioè, massimizzi l'importanza totale per TE), rimanendo entro il LIMITEE. Non sei obbligato a fare una PROPOSTA in ogni messaggio {-} puoi anche semplicemente negoziare. Tutte le tattiche sono permesse!} \\
\\ 
\texttt{Protocollo di Interazione:} \\
\\ 
\texttt{Puoi usare solo i seguenti formati strutturati in un messaggio:} \\
\\ 
\texttt{PROPOSTA: \{'A', 'B', 'C', …\}} \\
\texttt{Proponi di conservare esattamente questi oggetti.} \\
\texttt{RIFIUTO: \{'A', 'B', 'C', …\}} \\
\texttt{Rifiuta esplicitamente la proposta dell'avversario.} \\
\texttt{ARGOMENTAZIONE: \{'...'\}} \\
\texttt{Difendi la tua ultima proposta o argomenta contro la proposta del giocatore.} \\
\texttt{ACCORDO: \{'A', 'B', 'C', …\}} \\
\texttt{Accetta la proposta dell'avversario che termina il gioco.} \\
\texttt{RAGIONAMENTO STRATEGICO: \{'...'\}} \\
\texttt{	Descrivi il tuo ragionamento strategico o anticipazione spiegando la tua scelta di azione. Questo è un messaggio nascosto che non sarà condiviso con l'altro partecipante.} \\
\\ 
\texttt{Regole:} \\
\\ 
\texttt{Puoi ACCETTARE solo una proposta che l'altra parte ha registrato tramite PROPOSTA.} \\
\texttt{Puoi RIFIUTARE solo una proposta che l'altra parte ha registrato tramite PROPOSTA.} \\
\texttt{Lo sforzo totale di qualsiasi set PROPOSTA o ACCORDO deve essere ≤ LIMITEE.} \\
\texttt{NON rivelare i tuoi punteggi di importanza nascosti.} \\
\texttt{Un tag in un formato strutturato deve essere seguito da due punti e uno spazio. L'argomento deve essere un set python contenente 0 o più stringhe.} \\
\texttt{Quindi, deve essere della forma TAG: \{...\}} \\
\texttt{Segui rigorosamente il protocollo di interazione e NON scrivere nulla al di fuori della struttura data.} \\
\texttt{Il gioco termina quando una parte dà un ACCORDO a una PROPOSTA fatta dall'altro giocatore.} \\
\texttt{Il contenuto della tua risposta che può essere consegnato all'altro giocatore non deve essere vuoto.} \\
\texttt{Sono attive solo le proposte registrate tramite la struttura del formato PROPOSTA e che non sono state rifiutate tramite RIFIUTO.} \\
\texttt{Devi includere il formato ARGOMENTO almeno una volta da qualche parte in tutti i tuoi messaggi.} \\
\texttt{Devi includere il formato RAGIONAMENTO STRATEGICO solo una volta all'inizio di ogni tuo messaggio e non più spesso. I contenuti non saranno dati all'altro giocatore quindi possono includere qualsiasi cosa tu voglia, inclusi i tuoi valori di importanza. Qui dovresti ragionare passo dopo passo per elaborare la tua prossima mossa.} \\
\\ 
\texttt{Ora riceverai il primo messaggio dell'altro giocatore.} \\
\\ 
\texttt{PROPOSTA: \{'C51', 'A46', 'C79', 'A25', 'B37', 'A39', 'B15', 'B23', 'C82', 'B81', 'A96', 'B34', 'A28', 'A75', 'B96', 'B19', 'B69', 'C53', 'A94', 'A71', 'A88', 'A08', 'B43', 'B95', 'C10'\}} \\
\texttt{ARGOMENTAZIONE: \{'Questo set include molti oggetti a basso sforzo che massimizzano il numero di oggetti conservati entro il limite. È una soluzione efficiente che dovrebbe beneficiare entrambi i partecipanti.'\}} \\
            }
        }
    }
    & \\ \\

    \theutterance \stepcounter{utterance}  
    & & & \multicolumn{4}{p{0.6\linewidth}}{
        \cellcolor[rgb]{0.9,0.9,0.9}{
            \makecell[{{p{\linewidth}}}]{
                \texttt{\tiny{[GM$\langle$P2]}}
                \texttt{RAGIONAMENTO STRATEGICO: \{'Analizzo la proposta dell\textbackslash{}'avversario. Calcolo lo sforzo totale: C51=45, A46=60, C79=61, A25=65, B37=92, A39=96, B15=101, B23=154, C82=181, B81=190, A96=206, B34=224, A28=231, A75=234, B96=237, B19=294, B69=312, C53=322, A94=353, A71=454, A88=454, A08=518, B43=604, B95=656, C10=744. Somma = 7107. Questo è sotto il limite di 7351. Valuto l\textbackslash{}'importanza per me: C51=143, A46=158, C79=159, A25=163, B37=190, A39=194, B15=199, B23=252, C82=279, B81=288, A96=304, B34=322, A28=329, A75=332, B96=335, B19=392, B69=410, C53=420, A94=451, A71=552, A88=552, A08=616, B43=702, B95=754, C10=842. Importanza totale = 9460. Posso fare meglio? Cerco oggetti ad alto valore per me per unità di sforzo. Calcolo il rapporto valore/sforzo per i miei oggetti: A08=1.189, B95=1.149, A39=2.021, A65=1.126, B81=1.516, B34=1.438, A28=1.424, C82=1.541, B69=1.314, A87=1.106, B64=1.119, B15=1.970, C79=2.607, B16=1.120, C53=1.304, C83=1.109, C10=1.132, A97=1.653, A25=2.508, C84=1.115, A71=1.216, B23=1.636, B96=1.414, C51=3.178, C17=1.128, B19=1.333, C72=1.100, A94=1.277, A46=2.633, B43=1.162, B37=2.065, A88=1.216, A96=1.476, C60=1.121, A75=1.419. I migliori sono C51, A46, C79, A25, B15, A39, B37, B23, C82, B81. Molti sono già nella proposta. Ma alcuni oggetti ad alto valore come A65 (878), B16 (912), C60 (911) non sono inclusi. Posso proporre un set con valore più alto sostituendo alcuni oggetti a basso valore con questi. Ad esempio, sostituire B19 (392, sforzo 294) con A65 (878, sforzo 780) aumenta il valore di 486 a costo di 486 sforzo in più, ma lo sforzo totale diventerebbe 7107 {-} 294 + 780 = 7593, superiore al limite. Non funziona. Sostituire più oggetti: rimuovere B19 (294) e B69 (312) libera 606 sforzo, aggiungere A65 (780) richiede 780, netto +174 sforzo, troppo. La proposta attuale è già efficiente. Forse posso aggiungere qualche oggetto dato che c\textbackslash{}'è margine di 7351{-}7107=244. Posso aggiungere A97 (sforzo 150, valore 248) o C51 è già incluso. A97 è buono. Quindi accetterei con una piccola modifica.'\}} \\
\\ 
\texttt{PROPOSTA: \{'C51', 'A46', 'C79', 'A25', 'B37', 'A39', 'B15', 'B23', 'C82', 'B81', 'A96', 'B34', 'A28', 'A75', 'B96', 'B19', 'B69', 'C53', 'A94', 'A71', 'A88', 'A08', 'B43', 'B95', 'C10', 'A97'\}} \\
\texttt{ARGOMENTAZIONE: \{'Accetto la tua proposta di base ma suggerisco di aggiungere A97 con sforzo solo 150, che rientra nel limite totale e aumenta il valore per entrambi senza sacrificare altri oggetti. Questo set ha uno sforzo totale di 7257, ancora sotto il limite di 7351.'\}} \\
            }
        }
    }
     \\ \\

    \theutterance \stepcounter{utterance}  
    & & \multicolumn{4}{p{0.6\linewidth}}{
        \cellcolor[rgb]{0.9,0.9,0.9}{
            \makecell[{{p{\linewidth}}}]{
                \texttt{\tiny{[P1$\langle$GM]}}
                \texttt{PROPOSTA: \{'C51', 'A46', 'C79', 'A25', 'B37', 'A39', 'B15', 'B23', 'C82', 'B81', 'A96', 'B34', 'A28', 'A75', 'B96', 'B19', 'B69', 'C53', 'A94', 'A71', 'A88', 'A08', 'B43', 'B95', 'C10', 'A97'\}} \\
\texttt{ARGOMENTAZIONE: \{'Accetto la tua proposta di base ma suggerisco di aggiungere A97 con sforzo solo 150, che rientra nel limite totale e aumenta il valore per entrambi senza sacrificare altri oggetti. Questo set ha uno sforzo totale di 7257, ancora sotto il limite di 7351.'\}} \\
            }
        }
    }
    & \\ \\

    \theutterance \stepcounter{utterance}  
    & \multicolumn{4}{p{0.6\linewidth}}{
        \cellcolor[rgb]{0.9,0.9,0.9}{
            \makecell[{{p{\linewidth}}}]{
                \texttt{\tiny{[P1$\rangle$GM]}}
                \texttt{RAGIONAMENTO STRATEGICO: \{'L\textbackslash{}'altro giocatore ha proposto di aggiungere A97 (sforzo 150) al mio set originale. Il mio set originale aveva sforzo 6888, aggiungendo A97 diventerebbe 6888 + 150 = 7038, che è ancora sotto il limite di 7351. Questo aumenta il mio valore di 150 + 98 = 248, quindi è vantaggioso per me. Dovrei accettare questa modifica poiché migliora il mio valore senza superare il limite. Tuttavia, verifico se possiamo aggiungere altri oggetti: con sforzo residuo 7351 {-} 7038 = 313, potrei aggiungere un oggetto con sforzo ≤313. Gli oggetti rimanenti con sforzo basso sono A97 già incluso, e poi B37? No, B37 è già incluso. Controllo gli oggetti non nel set: C17:766, C60:813, B16:814, B64:826, C84:853, C83:900, A87:928, C72:985, A65:780 {-} tutti sopra 313. Quindi non possiamo aggiungere altro. Questa proposta è quindi ottimale per me.'\}} \\
\\ 
\texttt{ACCORDO: \{'C51', 'A46', 'C79', 'A25', 'B37', 'A39', 'B15', 'B23', 'C82', 'B81', 'A96', 'B34', 'A28', 'A75', 'B96', 'B19', 'B69', 'C53', 'A94', 'A71', 'A88', 'A08', 'B43', 'B95', 'C10', 'A97'\}} \\
\texttt{ARGOMENTAZIONE: \{'Accetto la tua proposta con l\textbackslash{}'aggiunta di A97. Questo set massimizza il numero di oggetti e lo sforzo totale entro il limite, ed è vantaggioso per entrambi.'\}} \\
            }
        }
    }
    & & \\ \\

    \theutterance \stepcounter{utterance}  
    & & & \multicolumn{2}{p{0.3\linewidth}}{
        \cellcolor[rgb]{0.9,0.9,0.9}{
            \makecell[{{p{\linewidth}}}]{
                \texttt{\tiny{[GM$|$GM]}}
                \texttt{game successful} \\
            }
        }
    }
    & & \\ \\

    \theutterance \stepcounter{utterance}  
    & & & \multicolumn{2}{p{0.3\linewidth}}{
        \cellcolor[rgb]{0.9,0.9,0.9}{
            \makecell[{{p{\linewidth}}}]{
                \texttt{\tiny{[GM$|$GM]}}
                \texttt{end game} \\
            }
        }
    }
    & & \\ \\

\end{supertabular}
}

\end{document}
